\clearpage
\section*{Appendix R\quad Computational Receipts}\label{app:receipts}
\addcontentsline{toc}{section}{Appendix R: Computational Receipts}
\small\noindent Each computation marked \rcptmarker\ in the text is verified by a runnable receipt in the \texttt{receipts/} directory that ships with this work; status \textsf{OK} = verified (traced, run, output matches the claim). This appendix is generated from the receipt index, not maintained by hand.\par\medskip
\begingroup\renewcommand{\arraystretch}{1.25}
\begin{description}
\item[\label{rcpt:P17_the_alpha_family_is_a_foliation}\texttt{P17\_the\_alpha\_family\_is\_a\_foliation}]\hfill\textsf{[OK]}\\
\textit{the substrate ontology / the c--Lambda null-ruling gauge} \ (THE ALPHA-FAMILY IS A FOLIATION OF THE AMBIENT BY THE LEVEL SETS OF ONE QUADRATIC FORM, WITH THE LIGHT CONE AS ITS SINGULAR MEMBER -- AND THE CONFOCAL ORTHOGONALITY THEOREM, WHICH A STANDING LEAD EXPECTED TO APPLY, IS VACUOUS HERE FOR A STRUCTURAL REASON. THE PREMISE IS CORRECT: in the EQUILATERAL case -- which is what this substrate is -- the confocal family IS the dilation family. ** BUT THE CLASSICAL THEOREM IS NOT 'CONFOCAL QUADRICS ARE ORTHOGONAL' FULL STOP: it is that through a generic point pass THREE members, one of each type, and THOSE meet pairwise orthogonally. The confocal equation is CUBIC in the parameter generically and LINEAR in the equilateral case -- asserted -- so exactly ONE member passes through any point and there is no second member for it to be orthogonal TO. The hypothesis fails, not the conclusion. ** The same fact without machinery: a point assigns ONE value to the form, so distinct-alpha substrates are DISJOINT and cannot meet at any angle. ** THE DEGENERATION, NAMED: equilateral is exactly the confocal family's degenerate case -- as the semi-axes come together the cubic loses two roots, the three types collapse, and the triply-orthogonal system collapses with them. So the theorem went unused for the BEST POSSIBLE REASON: this substrate sits at the one point of the confocal moduli space where it has no content. ** ** AND WHAT SURVIVES IS BETTER THAN THE THEOREM WOULD HAVE BEEN: the family is a foliation -- de Sitter above zero, THE LIGHT CONE at zero, the two-sheet hyperboloid below. Different-alpha substrates NEST, do not meet, and share ONE asymptote, which is the family's own degenerate member. **). 
\emph{Computes:} the equilateral reduction computed; the degree in the confocal parameter computed and asserted (3 generic, 1 equilateral); the three signature classes tabulated
 \emph{Bound:} ** AND THAT IS A REASON FOR THE c--Lambda NULL-RULING GAUGE'S alpha-INDEPENDENCE, FROM THE FAMILY RATHER THAN FROM THE GAUGE: the null structure is carried by the one leaf common to the whole family, so it cannot depend on which leaf one stands on. This paper already calls the cone 'the locked asymptote' without connecting it to the family it belongs to. AND ONE ROUTING THE LEAD DID NOT ASK FOR: the map between members is a DILATION, which is CONFORMAL and not an isometry, so the alpha-family is a CONFORMAL ORBIT and belongs in the conformal ledger rather than the quadric-geometry file where it was left -- an item filed under the machinery someone first reached for rather than the machinery that fits it **
\ \emph{Run:} \texttt{python3 receipts/P17\_geometric\_core\_paper/P17\_the\_alpha\_family\_is\_a\_foliation.py}
\item[\label{rcpt:P17_power_of_a_point}\texttt{P17\_power\_of\_a\_point}]\hfill\textsf{[OK]}\\
\textit{sec:power} \ (**!! c54.161: THE CIRCULARITY IS RETIRED** -- the power is now MEASURED from bisected secant products and compared against X0 taken from the substrate constraint, at 2 alpha and 3 alpha, so eq:powerisheight is no longer true by construction. **!! MARKED c54.157: THE CITING SENTENCE ATTRIBUTES TO THIS FILE FIGURES IT DOES NOT PRODUCE** -- the ds\textasciicircum{}2, the 600 random heights and the max 5.7e-14 all come from `P17\_power\_is\_null`. **And its own eq:powerisheight is circular: X0 is DEFINED as sqrt(\textbar{}X\textbar{}\textasciicircum{}2 - alpha\textasciicircum{}2), so 'power = X0\textasciicircum{}2' holds by construction.** What it does compute honestly is Euclid III.36 -- the secant invariance to 4.4e-16 over 500 sightlines. Original claim: power of a point pow(P)=\textbar{}X\textbar{}\textasciicircum{}2-alpha\textasciicircum{}2 (Steiner/Euclid III.36) = invariant: hinge power=3 exact; 500 random sightlines PX.PY=\textbar{}PO\textbar{}\textasciicircum{}2-alpha\textasciicircum{}2 to 4.4e-16). 
\emph{Computes:} every-sightline invariant test at machine precision
\ \emph{Run:} \texttt{python3 receipts/P17\_geometric\_core\_paper/P17\_power\_of\_a\_point.py}
\item[\label{rcpt:P17_power_is_null}\texttt{P17\_power\_is\_null}]\hfill\textsf{[OK]}\\
\textit{sec:power prop:tangentnull} \ (tangent to throat is NULL: tangent length=\textbar{}X0\textbar{}=sqrt(R\textasciicircum{}2-alpha\textasciicircum{}2) identity (Euclid=hyperboloid RHS); ds\textasciicircum{}2=-dX0\textasciicircum{}2+\textbar{}dX\textbar{}\textasciicircum{}2=0, max 5.68e-14 over 600 heights (=paper's 5.7e-14); power-of-a-point = null condition, same equation). 
\emph{Computes:} tangent-length identity + ds\textasciicircum{}2=0 at 600 random points
\ \emph{Run:} \texttt{python3 receipts/P17\_geometric\_core\_paper/P17\_power\_is\_null.py}
\item[\label{rcpt:P17_qm_S4_vs_S5}\texttt{P17\_qm\_S4\_vs\_S5}]\hfill\textsf{[OK]}\\
\textit{sec:ledger} \ (**!! CORRECTED c54.157: THE ROW CLAIMED THREE VERIFIED ITEMS AND ONE WAS VERIFIED.** The realified-generator antisymmetry is computed; **'minimal faithful real rep R\textasciicircum{}6' was a print string** (exhibiting one 6-dimensional rep shows su(3) FITS in so(6) and excludes nothing about 5), and 'GH beta = 2 pi alpha n-independent' has kappa hardcoded to 1.0 and never varies with n. **The minimality is now ENUMERATED by the Weyl dimension formula over su(3)'s irreps: the smallest faithful real dimension is 6, attained by 3 + 3-bar, so no faithful real rep fits in R\textasciicircum{}5.** Original claim: su(3) requires so(6)/S\textasciicircum{}5 not so(5)/S\textasciicircum{}4: 8 Gell-Mann generators realified to R\textasciicircum{}6 all antisymmetric (max\textbar{}R+R\textasciicircum{}T\textbar{}=0.0 -> so(6)), faithful on all 6 real dims (minimal real rep R\textasciicircum{}6); GH beta=2pi alpha n-independent -> thermal + gauge Euclideans co-locate on one real form (S\textasciicircum{}5 or S\textasciicircum{}4 equator)). 
\emph{Computes:} realified-generator antisymmetry + faithful-dimension + GH n-independence
\ \emph{Run:} \texttt{python3 receipts/P17\_geometric\_core\_paper/P17\_qm\_S4\_vs\_S5.py}
\item[\label{rcpt:C3_inversion_extends}\texttt{C3\_inversion\_extends}]\hfill\textsf{[OK]}\\
\textit{rem:onecircle} \ (C3 --- does INVERSION in the throat extend off the equatorial plane, as polarity did (Q1/Q6)? Minkowski inversion in the quadric eta(X,X)=alpha\textasciicircum{}2: P* = alpha\textasciicircum{}2 P / eta(P,P).). 
\emph{Computes:} runs clean; registered r2376 (c54) from its own docstring and a live run
 \emph{Bound:} description is the receipt's own wording, condensed; not independently re-derived here
\ \emph{Run:} \texttt{python3 receipts/P17\_geometric\_core\_paper/C3\_inversion\_extends.py}
\item[\label{rcpt:D1_hinge_duality}\texttt{D1\_hinge\_duality}]\hfill\textsf{[OK]}\\
\textit{rem:onecircle} \ (D1 --- PROJECTIVE DUALITY: the dual of the swinging door. PRIMAL : the pencil of slicing planes (doors) through the HINGE. DUALITY : in the plane of the throat, w.r.t. the throat circle \textbar{}x\textbar{}\textasciicircum{}2 = alpha\textasciicircum{}2, the pencil of LINES through a point P dualises to the RANGE OF POINTS on P's POLAR. So the dual of the swing is a POINT SLIDING ALONG THE HINGE'S POLAR. Where is that line?). 
\emph{Computes:} runs clean; registered r2376 (c54) from its own docstring and a live run
 \emph{Bound:} description is the receipt's own wording, condensed; not independently re-derived here
\ \emph{Run:} \texttt{python3 receipts/P17\_geometric\_core\_paper/D1\_hinge\_duality.py}
\item[\label{rcpt:D3_segre_no_pencil}\texttt{D3\_segre\_no\_pencil}]\hfill\textsf{[OK]}\\
\textit{rem:onecircle} \ (D3 --- SEGRE. The Segre symbol classifies a PENCIL of quadrics by the elementary divisors of lambda A + mu B. So the probe needs an honest answer to: DOES THE CORPUS HAVE TWO QUADRICS IN ONE SPACE? Confirmation 4, before anything runs. Candidates, each checked:). 
\emph{Computes:} runs clean; registered r2376 (c54) from its own docstring and a live run
 \emph{Bound:} description is the receipt's own wording, condensed; not independently re-derived here
\ \emph{Run:} \texttt{python3 receipts/P17\_geometric\_core\_paper/D3\_segre\_no\_pencil.py}
\item[\label{rcpt:O2_sightline_null_on_lift}\texttt{O2\_sightline\_null\_on\_lift}]\hfill\textsf{[OK]}\\
\textit{prop:tangentnull} \ (O2 --- p0's tangent/secant distinction, and whether it is a shadow boundary in the optics sense. p0: 'sightlines of the projection that TOUCH the throat are light rays; those that CUT it are not.'). 
\emph{Computes:} runs clean; registered r2376 (c54) from its own docstring and a live run
 \emph{Bound:} description is the receipt's own wording, condensed; not independently re-derived here
\ \emph{Run:} \texttt{python3 receipts/P17\_geometric\_core\_paper/O2\_sightline\_null\_on\_lift.py}
\item[\label{rcpt:Q1_quadric_polarity}\texttt{Q1\_quadric\_polarity}]\hfill\textsf{[OK]}\\
\textit{rem:onecircle} \ (Q1, invariant form. The first pass computed a EUCLIDEAN distance for the polar hyperplane in a MINKOWSKI space and got 1/sqrt7 = 0.37796 -- a gauge artifact, not an invariant. Redone with the only quantity polarity actually depends on: eta(P,P) against alpha\textasciicircum{}2.). 
\emph{Computes:} runs clean; registered r2376 (c54) from its own docstring and a live run
 \emph{Bound:} description is the receipt's own wording, condensed; not independently re-derived here
\ \emph{Run:} \texttt{python3 receipts/P17\_geometric\_core\_paper/Q1\_quadric\_polarity.py}
\item[\label{rcpt:Q3_cayley_klein}\texttt{Q3\_cayley\_klein}]\hfill\textsf{[OK]}\\
\textit{---} \ (Q3 --- Cayley-Klein: is the substrate's metric a log-cross-ratio against the absolute, and is alpha the CK scale constant? Substrate dS\_5 = \{X in M\textasciicircum{}6 : eta(X,X)=alpha\textasciicircum{}2\}.). 
\emph{Computes:} runs clean; registered r2376 (c54) from its own docstring and a live run
 \emph{Bound:} description is the receipt's own wording, condensed; not independently re-derived here
\ \emph{Run:} \texttt{python3 receipts/P17\_geometric\_core\_paper/Q3\_cayley\_klein.py}
\item[\label{rcpt:Q6r_polar_is_the_background}\texttt{Q6r\_polar\_is\_the\_background}]\hfill\textsf{[OK]}\\
\textit{---} \ (Q6-CONFIRMED --- the polar slice sits at RUNG 2, and that answers A18. CONFIRMATION 1 (ontology): section 0's five-rung ladder. Rung 1 = the substrate dS\_5 = SO(5,1)/SO(4,1). Rung 2 = 'A CUT of the substrate -- a four-geometry. dS\_4 is one such: the maximally symmetric one... A geometry is a cut exactly when its isometry group contains a sweep-subgroup of the substrate's.'). 
\emph{Computes:} runs clean; registered r2376 (c54) from its own docstring and a live run
 \emph{Bound:} description is the receipt's own wording, condensed; not independently re-derived here
\ \emph{Run:} \texttt{python3 receipts/P17\_geometric\_core\_paper/Q6r\_polar\_is\_the\_background.py}
\item[\label{rcpt:P17_no_second_scale_on_either_face}\texttt{P17\_no\_second\_scale\_on\_either\_face}]\hfill\textsf{[OK]}\\
\textit{the constant ledger} \ (**NEITHER REAL FORM OF THE SUBSTRATE SUPPLIES A SECOND SCALE, SO THE CONSTRUCTION CANNOT FORCE A DIMENSIONLESS MAGNITUDE. THE SILENCE IS A PROPERTY, NOT A GAP.** The Wick rotation X0 -> i X0 carries the defining quadric to S\textasciicircum{}5 of the SAME radius alpha -- it acts on the COORDINATE and leaves the right-hand side untouched, so no step exists at which a second radius could enter; the two faces are the two real forms of one complex quadric and a real form does not choose its own radius. Every curvature invariant on both faces is a pure power of 1/alpha\textasciicircum{}2 (asserted by logarithmic derivative), leaving no second constant with which to form a ratio. The remaining candidates are enumerated and each disposed: the horizon period is built from alpha; M and r\_0 are moduli; 2/(3 sqrt 3) is one member's value. **So a dimensionless magnitude cannot be forced, and that is the COMMON ROOT of three verdicts already banked separately -- the winding quantises and cannot measure, the flat bundle selects and cannot couple, the branch point filters and cannot supply.**). 
\emph{Computes:} the Wick image asserted equal to S\textasciicircum{}5 of radius alpha; each invariant's logarithmic derivative asserted an integer
 \emph{Bound:} **bound: this is about the two faces' GEOMETRY. Fields living on either face may carry their own scales -- but those are CONTENT, adjoined rather than derived, and adjoining one is what the corpus's admissibility criterion rejects. The escape exists and costs the thing the programme is for**
\ \emph{Run:} \texttt{python3 receipts/P17\_geometric\_core\_paper/P17\_no\_second\_scale\_on\_either\_face.py}
\item[\label{rcpt:P03_operator_at_general_D}\texttt{P03\_operator\_at\_general\_D}]\hfill\textsf{[OK]}\\
\textit{rem:dimension / thm:kernel} \ (** ASSUMPTION (A1) DISCHARGED. ** CR's OWN operator at general D returns Tangherlini-de Sitter and nothing else. The point previously missed: CR's operator is NOT a metric function -- P8 thm:kernel defines the vacuum sector as the KERNEL OF THE MATTER FUNCTIONAL on a cut in the construction gauge ('derived, not matched'), and that condition runs at any D. Computed from the metric rather than quoted: for ds\textasciicircum{}2 = -f dt\textasciicircum{}2 + dr\textasciicircum{}2/f + r\textasciicircum{}2 dOmega\textasciicircum{}2\_\{D-2\}, G\textasciicircum{}t\_t = (D-2)/(2r\textasciicircum{}2)[r f' + (D-3)(f-1)] (verified at D=4,5,6), so T=0 is the first-order linear ODE r f' + (D-3)(f-1) + 2 Lambda r\textasciicircum{}2/(D-2) = 0 -- P8's own equation at D=4 -- whose ENTIRE solution space at D=4,5,6,7,8 is f = 1 - 2M/r\textasciicircum{}(D-3) - 2 Lambda r\textasciicircum{}2/((D-1)(D-2)) with ONE constant of integration, verified identical to Tangherlini-dS at every D tested. With alpha(D) = sqrt((D-1)(D-2)/2Lambda), which returns sqrt(3/Lambda) at four, this is f = 1 - 2M/r\textasciicircum{}(D-3) - r\textasciicircum{}2/alpha\textasciicircum{}2 at every D, and the dimension result's input 2M = r\_0\textasciicircum{}(D-3) - r\_0\textasciicircum{}(D-1) is that f at a root. THE DIMENSION RESULT THEREFORE STRENGTHENS: the collapse, the fold being D-1 and the mass parity all run on a function the construction GENERATES rather than borrows [borrowed from P3 / P8]). 
\emph{Computes:} the Einstein tensor computed from the metric at D=4,5,6 and matched to the closed form; the ODE solved by dsolve at D=4..8 and compared term by term with Tangherlini-dS
 \emph{Bound:} TWO INPUTS REMAIN AND ARE NAMED: (i) THE LOCK -- the construction gauge g\_tt g\_rr = -1 is assumed at general D, and P8 says the single slicing curve enforces it at four but this does not show it does so at general D (L-89; P8's lapse split is the machinery that would say otherwise); (ii) THE SPHERE -- the transverse space is taken round S\textasciicircum{}\{D-2\}, which at general D is a choice (L-90). And (A2), the single-harmonic collapse being a REQUIREMENT at other D, is untouched
\ \emph{Run:} \texttt{python3 receipts/P03\_SdS\_slicing/P03\_operator\_at\_general\_D.py}
\item[\label{rcpt:P17_the_cut_is_planar_only_at_M_zero}\texttt{P17\_the\_cut\_is\_planar\_only\_at\_M\_zero}]\hfill\textsf{[OK]}\\
\textit{the ladder / sec:slicing / sec:pd} \ (**!! c54.161: THE TWO c54.157 MARKINGS ARE RETIRED** -- the hyperplane split is now done for a tilted n with the induced Ricci scalar computed as 12/(alpha\textasciicircum{}2-c\textasciicircum{}2), and the Kretschmann scalar is DERIVED from the metric rather than typed in. **!! CORRECTED c54.157: THE ROW'S 'the hyperplane split done in general' IS FOUR PRINT STATEMENTS AND A HARDCODED TABLE**, and the Kretschmann scalar 48M\textasciicircum{}2/r\textasciicircum{}6 + 24/alpha\textasciicircum{}4 is an INPUT typed at line 97, not a finding -- it is derived from the metric in `P07\_sds\_kretschmann`, which is where the paper should look. The remaining evidence items (the M-factored closure defect, the r-dependent kappa\textasciicircum{}2, the M = 0 control) are accurate. Original claim: WHAT "CUT" CAN MEAN, COMPUTED. EVERY HYPERPLANE SECTION OF dS\_5 HAS CONSTANT CURVATURE -- splitting X = c n + Y gives eta(Y,Y) = alpha\textasciicircum{}2 - c\textasciicircum{}2, a dS\_4 of radius sqrt(alpha\textasciicircum{}2-c\textasciicircum{}2), for every n and c; a quadric cut by a linear space is a quadric. And SdS has K = 48 M\textasciicircum{}2/r\textasciicircum{}6 + 24/alpha\textasciicircum{}4, non-constant unless M = 0. ** SO A PLANE SECTION OF dS\_5 IS SdS IF AND ONLY IF M = 0 ** -- which CONFIRMS p0's polar hyperplane and shows it is not one example but THE only one. HENCE MATTER MUST BEND THE CUT, now a theorem rather than a reading. AND THE BENT CUT DOES NOT CLOSE EITHER: with P3's own symmetry (round areal sphere in a spacelike 3-block, staticity as the ambient boost the pencil turns on), lying ON the substrate fixes G and inducing SdS fixes G' -- two determinations of one function, and solving for kappa\textasciicircum{}2 returns an r-DEPENDENT expression, contradicting kappa's definition. ** AND THE FORM OF THE OBSTRUCTION IS THE RESULT: the closure defect factors with an OVERALL FACTOR OF M. The failure of a cut to close as a hypersurface IS the mass it carries, linearly, with no other source ** [borrowed from P17 / P3 / P9]). 
\emph{Computes:} the hyperplane split done in general; the Kretschmann derivative; the two determinations of G' computed and subtracted; kappa\textasciicircum{}2 solved for and shown r-dependent; the M = 0 control returning residual 0 and G = 0 (the planar vacuum cut)
 \emph{Bound:} ** IT DOES NOT REFUTE THE CUT RELATION, and reading it that way would be the error the receipt exists to prevent: P9 defines a cut by its isometry group containing a sweep-subgroup, P12 by an orientation plus causal-vantage datum, and the framework calls SdS a DERIVATIVE metric, not an induced one. What is established is that the ISOMETRIC-EMBEDDING reading is available at M=0 and nowhere else, so the corpus owes a sentence saying which relation 'cut' names **
\ \emph{Run:} \texttt{python3 receipts/P17\_geometric\_core\_paper/P17\_the\_cut\_is\_planar\_only\_at\_M\_zero.py}
\item[\label{rcpt:P0_the_order_parameter_is_the_offset_and_it_is_bounded_by_the_nariai_member}\texttt{P0\_the\_order\_parameter\_is\_the\_offset\_and\_it\_is\_bounded\_by\_the\_nariai\_member}]\hfill\textsf{[OK]}\\
\textit{sec:residue} \ (**ITEM 48 WORKED: THE TWO SYMMETRY BREAKINGS ARE ONE OBJECT, AND JOINING THEM GIVES THE HIGGS IDENTIFICATION ACTUAL CONTENT.** Lead `L-521`; VEIN `L-221`. **PART 1** p0 states the parity as \$r\_0\textbackslash\{\}mapsto-r\_0\$ (whence \$2M\textbackslash\{\}mapsto-2M\$); P3 derives \$2M=r\_0-r\_0\textasciicircum{}3\$ for a different purpose; **the cubic is odd, so p0's mass-reflection is an IDENTITY of P3's relation** and the two breakings are one object. **PART 2** \$2M=0\$ at \$r\_0=0,\textbackslash\{\}pm\textbackslash\{\}alpha\$ with \$f=1-r\textasciicircum{}2/\textbackslash\{\}alpha\textasciicircum{}2\$ at all three --- **the symmetric sector is the bare substrate**, one geometry read at three offsets. **PART 3** over P3's admissible range \$\textbackslash\{\}lvert r\_0\textbackslash\{\}rvert\textbackslash\{\}le2\textbackslash\{\}alpha/\textbackslash\{\}sqrt3\$, **\$\textbackslash\{\}lvert M\textbackslash\{\}rvert\textbackslash\{\}le\textbackslash\{\}alpha/3\textbackslash\{\}sqrt3\$ --- the Nariai mass, derived independently from \$f=f'=0\$**; a quartic potential is unbounded in its field and this order parameter is not, saturating at the configuration a collapse reaches. **PART 4** eight source checks across p0 and P3, including that p0 now names the Higgs (zero uses in seventeen papers) and keeps `F1` explicit. [borrowed from p0]). 
\emph{Computes:} **the oddness of \$2M(r\_0)\$ --- this IS the join; if P3's relation were not odd, p0's mass-reflection would be a separate assumption and there would be no single object**; all three roots asserted to give the same de Sitter metric function, both as the symmetric-sector claim and so "three roots" cannot be read as three vacua; the bound asserted numerically over P3's own range AND identified with the Nariai mass **derived independently and never substituted in, because the whole content is that two separately-derived numbers coincide**; and eight source checks
 \emph{Bound:} **scope: a statement about the MECHANISM, which is what item 48 asked for. NOT claimed: any vacuum expectation value, electroweak scale or mass value; any derivation of the Higgs field, its potential or its gauge group. `F1` untouched and said so in the paper's own voice. And the three roots are three offsets on ONE geometry --- vantage multiplicity already in the corpus.**
\ \emph{Run:} \texttt{python3 receipts/P17\_geometric\_core\_paper/P0\_the\_order\_parameter\_is\_the\_offset\_and\_it\_is\_bounded\_by\_the\_nariai\_member.py}
\end{description}\endgroup
